% The "cherry tree": each branch is a rational prime (p), and the
% cherries below it are the finitely many prime ideals dividing (p).
% Source: book/tex/alg-NT/classgrp.tex (§ The class group is finite) —
% asy figure.
\documentclass[border=2mm]{standalone}
\usepackage{tikz}
\begin{document}
\begin{tikzpicture}[scale=1.4]
  \node[above] at (1,0) {$(2)$};
  \foreach \i in {1,...,4} \fill (1,-\i/4) circle (1.4pt);
  \node[above] at (2,0) {$(3)$};
  \foreach \i in {1,...,6} \fill (2,-\i/4) circle (1.4pt);
  \node[above] at (3,0) {$(5)$};
  \foreach \i in {1,...,2} \fill (3,-\i/4) circle (1.4pt);
  \node[above] at (4,0) {$\dots$};
\end{tikzpicture}
\end{document}
